%% ----------------------------------------------------------------
%% MCM/ICM 2025 Template - Fixed for VS Code Local Environment
%% ----------------------------------------------------------------

% 【修改1】demo模式仅用于图片占位，提交前删除[demo]
\documentclass[demo]{mcmthesis}

\mcmsetup{
  CTeX = false,           % 英文论文关闭CTeX
  tcn = 2602787,          % 团队编号（核心，务必核对）
  problem = B,            % 题目类型
  sheet = true,           
  titleinsheet = true, 
  keywordsinsheet = true, 
  titlepage = false,
  abstract = true
}

\usepackage{palatino} 
\usepackage{lipsum}   
\usepackage{amsmath, amssymb, amsfonts} 
\usepackage{booktabs} 
\usepackage{geometry} 
\usepackage{float}    
\usepackage{indentfirst}
\usepackage{subcaption} 
\usepackage{graphicx}
\usepackage{fancyhdr}
\usepackage{lastpage}
\usepackage{silence}
\usepackage{arydshln} % for \cdashline

\WarningFilter{everypage}{}%忽略everypage警告

% 【修改点2】删除了 \usepackage[demo]{graphicx}，避免冲突

% 【修改点3】解决 "headheight is too small" 的黄色警告
\setlength{\headheight}{15pt}

% 【修改5】页码样式：添加“当前页/总页数”，符合美赛习惯
\pagestyle{fancy}

\cfoot{\thepage/\pageref{LastPage}} % 页脚居中显示“页码/总页数”

% 备忘录背景（保留原有逻辑，修复路径问题）（可选）
\usepackage{background}
\backgroundsetup{
    scale=1,
    angle=0,
    opacity=0.1,  
    contents={}   % 默认无背景（正文页）
}

% --- 标题信息 ---
\title{Your Title Here}
\author{Team \#2602787}
\date{\today}

% ================= 文档开始 =================
\begin{document}

% --- 摘要页（美赛要求单独一页）---
\begin{abstract}
    The development of tourism can lead to economic growth, but excessive development can cause ecological and social issues. This paper analyzes the situation of tourism in Juneau, Alaska, and formulates an optimization model...

    For Problem 1, ...

    For Problem 2, ...

    For Problem 3...

    \begin{keywords}
        Sustainable tourism; Optimization; MCM
    \end{keywords}
\end{abstract}

\maketitle

% --- 目录（单独一页，美赛要求）---
\tableofcontents
\newpage

% ================= 第一章 =================
\section{Introduction}
\subsection{Problem Background}
\setlength{\parindent}{2em}
As humanity's technological capabilities continue to advance, the demand for scalable, safe, and environmentally friendly access to space has become increasingly critical. Traditional rocket-based launch systems have enabled early space exploration; however, their high costs, limited payload capacity, and environmental impacts pose significant challenges to long-term extraterrestrial development.

On one hand, emerging technologies such as space elevator systems offer a promising alternative by enabling continuous, electrically powered transportation of materials from Earth to orbit, potentially reducing costs and atmospheric pollution. On the other hand, the construction of a large-scale lunar colony requires the delivery of enormous quantities of materials and supplies, placing stringent demands on transportation capacity, reliability, and timelines.

How to balance cost, efficiency, and feasibility among different transportation strategies is therefore central to the successful construction and sustainable operation of a 100,000-person lunar colony. Addressing this challenge requires a quantitative evaluation of alternative logistics schemes through mathematical modeling and systematic comparison.


\subsection{Restatement of the Problem}

Based on the above background, this study aims to:
\begin{itemize}
    \item \textbf{Develop a Lunar Transportation Evaluation Model:} 
    Formulate a mathematical framework to describe large-scale material transportation from Earth to the Moon, and compare different strategies including the Space Elevator System, traditional rocket launches, and their combinations in terms of cost, timeline, and feasibility.

    \item \textbf{Analyze System Robustness and Supply Requirements:} 
    Examine how non-ideal operating conditions, such as reduced system efficiency or transportation failures, affect model outcomes, and extend the model to evaluate the additional cost and timeline required for long-term water and supply delivery.

    \item \textbf{Assess Environmental Impact and Provide Strategic Recommendations:} 
    Investigate the environmental implications of different transportation strategies and summarize the results in a decision-oriented memorandum recommending an effective and sustainable course of action for constructing and sustaining a 100,000-person Moon Colony.
\end{itemize}


\subsection{Our Work}
To avoid complicated description...
% 【插入流程图】
% 因为开启了 demo 模式，这里会显示一个黑框，不会报错找不到图片
\begin{figure}[H]
    \centering
    \includegraphics[width=0.9\textwidth]{flowchart.jpg}
    \caption{Overview of our work}
    \label{fig:flowchart}
\end{figure}

% ================= 第二章 =================
\section{Model Preparations}

\subsection{Assumptions and Justifications}
\begin{itemize}
    \item \textbf{Assumption 1:} Assumption text here...
    \item \textbf{Justification:} Justification text here...

    \item \textbf{Assumption 2:} Assumption text here...
    \item \textbf{Justification:} Justification text here...
\end{itemize}

\subsection{Notations}
\renewcommand{\arraystretch}{1.25}
\setlength{\tabcolsep}{10pt}

\begin{longtable}{>{\centering\arraybackslash}p{0.10\linewidth}
    >{\centering\arraybackslash}p{0.60\linewidth}
    >{\centering\arraybackslash}p{0.16\linewidth}}

    \caption{Symbol Definitions and Notation}\label{tab:notations}                                                                                                  \\

    \toprule
    \textit{Symbols}             & \textit{Description}                                                                                    & \textit{Unit}          \\
    \midrule
    \endfirsthead

    \toprule
    \textit{Symbols}             & \textit{Description}                                                                                    & \textit{Unit}          \\
    \midrule
    \endhead

    \midrule
    \multicolumn{3}{r}{\textit{Continued on next page}}                                                                                                             \\
    \endfoot

    \bottomrule
    \endlastfoot

    % ======= your rows start =======

    \multicolumn{3}{c}{\textit{\textbf{Sets and Indices}}}                                                                                                          \\
    \cdashline{1-3}

    $i$                          & Transportation method index ($i \in \{\text{SE}, \text{RK}\}$, where SE = Space Elevator, RK = Rockets) & --                     \\
    $j$                          & Galactic Harbour index ($j = 1,2,3$)                                                                    & --                     \\
    $k$                          & Rocket launch site index ($k = 1,2,3,\dots,10$)                                                         & --                     \\

    \cdashline{1-3}
    & & \\
    \multicolumn{3}{c}{\textit{\textbf{Global Parameters}}}                                                                                                         \\
    \cdashline{1-3}

    $t$                          & Time (years), $t \ge 2050$                                                                              & year                   \\
    $M_{\text{total}}$           & Total materials required for colony construction                                                        & metric ton             \\
    $M_{\text{colony}}$          & Designed population capacity of the Moon Colony                                                         & person                 \\
    $T_0$                        & Start year of transportation (2025)                                                                     & year                   \\
    $x$                          & Fraction transported by Space Elevator $(0\le x\le 1)$                                                  & --                     \\

    \cdashline{1-3}
    \multicolumn{3}{c}{\textit{\textbf{Space Elevator System Parameters}}}                                                                                          \\
    \cdashline{1-3}

    $C_{\text{SE}}$              & Cost per metric ton via Space Elevator                                                                  & USD/ton                \\
    $Q_{\text{SE}}$              & Annual capacity of one Galactic Harbour                                                                 & ton/year               \\
    $Q_{\text{SE}}^{\text{tot}}$ & Total annual capacity of three Galactic Harbours                                                        & ton/year               \\
    $\eta_{\text{SE}}$           & Operational efficiency of Space Elevator System                                                         & --                     \\
    $T_{\text{SE}}$              & Time to deliver required mass via Space Elevator                                                        & year                   \\
    $M_{\text{SE}}$              & Mass delivered by Space Elevator ($xM_{\text{total}}$)                                                  & ton                    \\

    \cdashline{1-3}
    \multicolumn{3}{c}{\textit{\textbf{Rocket Transportation Parameters}}}                                                                                          \\
    \cdashline{1-3}

    $C_{\text{RK}}$              & Cost per metric ton via rockets                                                                         & USD/ton                \\
    $P_{\text{RK}}$              & Payload per rocket launch                                                                               & ton                    \\
    $N_{\text{RK}}$              & Number of launches per year                                                                             & 1/year                 \\
    $\eta_{\text{RK}}$           & Rocket launch success rate                                                                              & --                     \\
    $T_{\text{RK}}$              & Time to deliver required mass via rockets                                                               & year                   \\
    $M_{\text{RK}}$              & Mass delivered by rockets ($(1-x)M_{\text{total}}$)                                                     & ton                    \\

    \cdashline{1-3}
    \multicolumn{3}{c}{\textit{\textbf{Water and Sustenance Notations}}}                                                                                            \\
    \cdashline{1-3}

    $W_d$                        & Daily water demand per person                                                                           & ton/(person$\cdot$day) \\
    $W_{\text{year}}$            & Annual water requirement for the colony                                                                 & ton                    \\
    $M_{\text{sup}}$             & Annual mass of routine supplies                                                                         & ton                    \\

    \cdashline{1-3}
    \multicolumn{3}{c}{\textit{\textbf{Cost and Performance Metrics}}}                                                                                              \\
    \cdashline{1-3}

    $\text{Cost}_{\text{SE}}$    & Total cost of Space Elevator transportation                                                             & USD                    \\
    $\text{Cost}_{\text{RK}}$    & Total cost of rocket transportation                                                                     & USD                    \\
    $\text{Cost}_{\text{total}}$ & Total transportation cost                                                                               & USD                    \\
    $T_{\text{finish}}$          & Total completion time of construction delivery                                                          & year                   \\

    \cdashline{1-3}
    \multicolumn{3}{c}{\textit{\textbf{Environmental Impact Indicators}}}                                                                                           \\
    \cdashline{1-3}

    $E_{\text{SE}}$              & Environmental impact index per ton (Space Elevator)                                                     & --                     \\
    $E_{\text{RK}}$              & Environmental impact index per ton (Rockets)                                                            & --                     \\
    $E_{\text{total}}$           & Total environmental impact index                                                                        & --                     \\

    % ======= your rows end =======
\end{longtable}



\subsection{Basic components}
The basic components... (Figure 3)

% ================= 第三章：模型设计 =================
\section{Model Design}

\subsection{Model Overview}
In this section, we develop mathematical models to evaluate different transportation strategies for delivering large-scale materials from Earth to the Moon. 
Three scenarios are considered: transportation solely via the Space Elevator System, transportation solely via traditional rocket launches, and a hybrid strategy combining both methods.

For each scenario, we analyze the total transportation cost, required completion time, and overall feasibility under ideal operating conditions. 
These baseline models provide a unified framework for subsequent sensitivity and robustness analyses.

\subsection{Space Elevator Transportation Model}
The Space Elevator System enables continuous and electrically powered transportation of materials from Earth to geosynchronous orbit through multiple Galactic Harbors. 
Let $Q_{\text{SE}}$ denote the annual transportation capacity of a single Galactic Harbor, and $C_{\text{SE}}$ represent the unit transportation cost per metric ton.

Assuming stable operation and constant efficiency, the total annual capacity of the Space Elevator System is given by
\begin{equation}
Q_{\text{SE}}^{\text{tot}} = 3 Q_{\text{SE}}.
\end{equation}

The total time required to deliver the necessary construction materials $M_{\text{total}}$ via the Space Elevator System can be expressed as
\begin{equation}
T_{\text{SE}} = \frac{M_{\text{total}}}{\eta_{\text{SE}} Q_{\text{SE}}^{\text{tot}}},
\end{equation}
where $\eta_{\text{SE}}$ denotes the operational efficiency of the system.

The corresponding total transportation cost is calculated as
\begin{equation}
\text{Cost}_{\text{SE}} = C_{\text{SE}} \cdot M_{\text{total}}.
\end{equation}

\subsection{Rocket Transportation Cost Model}

In addition to the Space Elevator System, traditional rocket launches are considered as an alternative means of transporting construction materials and supplies to the Moon. 
Due to rapid technological development and accumulated launch experience, the unit cost of rocket transportation is expected to decrease over time. 
To capture this trend, we adopt a learning-curve-based cost model inspired by Wright’s Law, which has been widely applied in industrial production and aerospace engineering.

The core assumption of Wright’s Law is that each doubling of cumulative production or operational experience leads to a fixed percentage reduction in unit cost. 
Let $C(t)$ denote the unit cost of a single rocket transportation mission at time $t$, $C_0$ represent the initial launch cost, and $N(t)$ be the cumulative number of rocket launches completed by time $t$. 
The cost model is expressed as
\begin{equation}
C(t) = C_0 \cdot [N(t)]^{\alpha(t)},
\end{equation}
where $\alpha(t)$ is the learning rate, reflecting technological improvement and operational efficiency gains.

Based on publicly available data from SpaceX and related industry reports, the initial cost of a heavy-lift rocket mission is estimated to be approximately \$450 million. 
Considering that learning effects differ across subsystems—such as propulsion, manufacturing, labor, and maintenance—we adopt an initial learning rate of $\alpha = 15\%$, which lies within the range reported in existing literature.

Furthermore, as technological maturity increases, the marginal benefit from learning is expected to diminish. 
To reflect this effect, the learning rate is modeled as a time-dependent function:
\begin{equation}
\alpha(t) = -0.1413 \cdot e^{\ln(2)\, t / 24} + 0.2913,
\end{equation}
which is derived through literature review, curve fitting, and extrapolation.

To estimate the cumulative number of rocket launches, historical launch data of SpaceX Starship missions are collected and analyzed. 
At early stages, the launch frequency grows rapidly due to technological breakthroughs, while in later stages it gradually stabilizes. 
Based on this observation, the cumulative launch function is approximated as
\begin{equation}
N(t) = k t + m,
\end{equation}
where $k$ and $m$ are parameters obtained through regression analysis.

Combining the above components, the unit cost of a single rocket transportation mission over time can be determined. 
Using this model, we estimate that by the year 2050, the cost of a single rocket transportation mission will decrease to approximately \$75 million. 
For computational convenience, this value is adopted as the representative unit cost for rocket-based transportation in subsequent analyses.


\subsection{Hybrid Transportation Strategy}
In practical applications, a hybrid transportation strategy combining the Space Elevator System and traditional rocket launches may offer improved flexibility and resilience.

Let $x$ denote the fraction of total materials transported via the Space Elevator System. 
Accordingly, the remaining fraction $(1-x)$ is delivered by rocket launches.
The total transportation cost under the hybrid strategy is expressed as
\begin{equation}
\text{Cost}_{\text{total}} = x \cdot \text{Cost}_{\text{SE}} + (1-x) \cdot \text{Cost}_{\text{RK}}.
\end{equation}

The overall completion time is determined by the slower of the two transportation processes.

\subsection{Solution Method and Comparison}
For each transportation scenario, we compute the total cost and completion time using the models established above. 
By comparing these performance indicators across the Space Elevator, rocket-only, and hybrid strategies, we evaluate their relative advantages and limitations in supporting the construction of a 100,000-person Moon Colony.

The results obtained in this section serve as the baseline for subsequent sensitivity analysis and robustness testing under non-ideal operating conditions.

% ================= 第4章：敏感性分析 =================
\section{Sensitivity Analysis}

\subsection{Parameter Sensitivity}
% 参数敏感性分析

\subsection{Robustness Testing}
% 鲁棒性测试

% ================= 第5章：模型应用 =================
\section{Application of Our Model}

\subsection{Case Study}
% 案例研究

\subsection{Results and Discussion}
% 结果与讨论

% ================= 第六章 =================
\section{Model Evaluation}

\subsection{Strengths}
\begin{itemize}
    \item \textbf{Strength 1:} ...
    \item \textbf{Strength 2:} ...
\end{itemize}

\subsection{Weaknesses and Future Improvements}
\begin{itemize}
    \item \textbf{Weakness 1:} ...
    \item \textbf{Weakness 2:} ...
\end{itemize}

% ================= 参考文献（美赛要求规范格式）=================
\newpage
\begin{thebibliography}{99}

\bibitem{1}
Lam, K. C., Lee, D., \& Hu, T. (2001).
Understanding the effect of the learning--forgetting phenomenon on project duration.
\textit{International Journal of Project Management}, 19(7), 411--420.
https://doi.org/10.1016/S0263-7863(00)00025-9

\bibitem{2}
Kim, S., Koo, J., \& Yoon, E. S. (2012).
Optimization of sustainable energy planning with consideration of uncertainties in learning rates and external cost factors.
In \textit{Computer Aided Chemical Engineering} (Vol. 31, pp. 871--875).
Elsevier.
https://doi.org/10.1016/B978-0-444-54298-4.50161-6

\bibitem{3}
Swan, P. (2019, August 18).
Myths busted: The real reason launch costs are high and how space access infrastructure can reduce launch costs to LEO.
Conference presentation at the International Space Elevator Conference.

\bibitem{4}
SpaceX. (n.d.).
\textit{Starship}.
Retrieved January 30, 2026, from https://www.spacex.com/vehicles/starship

\bibitem{5}
Young, W. A., II, Masel, D. T., \& Judd, R. P. (2007).
A matrix-based methodology for determining a part family’s learning rate.
\textit{Computers \& Industrial Engineering}, 53(4), 803--811.
https://doi.org/10.1016/j.cie.2007.08.002

\end{thebibliography}


% ================= 备忘录 (最后一页) =================
\newpage
% 开启背景图片
% 注意：即使是 demo 模式，最好也确保目录下有个随便叫什么名字的jpg文件，或者先把这行注释掉
\backgroundsetup{
    contents={\includegraphics[width=\paperwidth,height=\paperheight,keepaspectratio]{snow_mountain_bg.jpg}},
    scale=1,
    angle=0,
    opacity=0.2
}
\thispagestyle{empty}

\begin{center}
    {\Huge \textbf{MEMORANDUM}}
\end{center}
\vspace{1cm}

\large
\textbf{To:} The Tourist Council of Juneau \\
\textbf{From:} Team \#2602787 \\
\textbf{Subject:} Policy Advice on the Sustainable Tourism of Juneau \\
\textbf{Date:} January 27, 2025
\hrule
\vspace{1cm}

\normalsize
% 备忘录正文内容


\vspace{1cm}
Sincerely,

Team \#2602787

\end{document}